\chapter{Introdução}

\section{Contextualização e motivação}

A transparência e o acesso à informação configuram-se como demandas sociais em constante expansão e têm assumido papel central nos debates na esfera pública. Parte-se do pressuposto de que a adoção de práticas voltadas à promoção da transparência e ao combate à corrupção gera benefícios sistêmicos para a sociedade, ao fortalecer as instituições e assegurar que os recursos públicos sejam alocados de forma alinhada às necessidades prioritárias \cite{figueiredo2016papel}. 

Em resposta a essas reivindicações, o Poder Legislativo brasileiro promulgou a Lei de Acesso à Informação (LAI), a qual consolida o princípio de que o acesso às informações institucionais deve constituir a regra, e não a exceção \cite{L12527}. Nesse contexto, a transparência também passa a ser compreendida como um fator estratégico para a inovação e o desenvolvimento organizacional.

Os projetos constituem o principal vetor de inovação e transformação nos negócios; efetivamente, a única forma de modificar uma organização, implementar uma estratégia, inovar ou obter vantagem competitiva é por meio de projetos \cite{van20147}. Assegurar visibilidade a esses projetos é fundamental para o fortalecimento de uma organização, uma vez que, além de possibilitar o compartilhamento de conhecimento, viabiliza a atração de investimentos econômicos e a geração de impacto social.

A colaboração e o intercâmbio de conhecimento entre organizações \cite{nooshinfard2014success} e entre projetos \cite{wiewiora2014interactions} configuram uma realidade consolidada. Nesse cenário, os projetos precisam estar preparados para aproveitar os benefícios decorrentes dessa abertura de fronteiras e da circulação ampliada de informações \cite{vieira2020universal}.

A transparência é uma característica universal do desenvolvimento de projetos e dos processos de implementação. No entanto, as formas de operacionalizá-la variam conforme a metodologia adotada. No modelo clássico de gerenciamento de projetos, a transparência se restringe principalmente à função informativa, voltada à comunicação com as partes interessadas por meio de práticas tradicionais, como reuniões, avisos e relatórios formais.

Já modelos mais contemporâneos, que elevam a transparência ao nível axiológico, ampliam essa noção ao incorporar ferramentas que favorecem a compreensão dos processos e procedimentos de gestão. Essas abordagens promovem um fluxo bidirecional de informações e fortalecem a confiança pessoal e institucional entre os participantes do projeto \cite{betta2018transparency}. Nesse sentido, mecanismos estruturados para garantir essa transparência tornam-se essenciais para potencializar seus benefícios.

Um observatório de projetos desempenha papel central no fortalecimento da transparência ao promover a centralização, análise e disseminação sistemática de informações relacionadas aos projetos. Concebido como uma entidade dedicada à observação contínua — seja na forma de uma unidade organizacional, de um mecanismo ou processo estruturado, ou ainda de um sistema de informação —, o observatório estabelece um canal institucionalizado para o compartilhamento de dados relevantes.

A disponibilização de métricas, a visibilidade dos processos e a oferta de análises detalhadas ampliam a compreensão sobre o andamento dos projetos, favorecendo a comunicação entre as partes interessadas e sustentando práticas efetivas de prestação de contas. Além disso, ao estimular a colaboração, apoiar estudos e pesquisas e fornecer subsídios informacionais consistentes, os observatórios criam uma base sólida para a inovação, contribuindo para uma gestão de projetos mais eficiente, confiável e orientada à tomada de decisão informada \cite{vieira2021model}.

Nesse contexto, o Modelo para Observatório de Projetos (MPO) surge como uma referência estruturada para a compreensão e o desenvolvimento de observatórios de projetos. Proposto por \citeauthor{vieira2022thesis}, o MPO organiza 61 conceitos distribuídos em três níveis hierárquicos — geral, intermediário e específico — que descrevem de forma sistemática as principais estruturas, processos e agentes envolvidos nesses ambientes.

Ao sistematizar esses elementos, o modelo representa um avanço relevante no campo, ao oferecer uma base conceitual abrangente, empiricamente validada e capaz de orientar a implementação de iniciativas concretas de observação de projetos em diferentes contextos organizacionais \cite{vieira2022thesis}. 

Apesar dessa contribuição, o MPO encontra-se descrito predominantemente em linguagem natural e em estruturas conceituais, o que pode limitar sua precisão semântica, dificultar sua interpretação computacional e restringir sua utilização em sistemas de informação. Dessa forma, observa-se a necessidade de uma formalização ontológica do modelo, capaz de explicitar de maneira rigorosa seus conceitos, relações e restrições semânticas, permitindo sua representação formal e ampliando seu potencial de aplicação em ambientes computacionais.

\section{Descrição do problema}
Apesar de sua robustez conceitual e validação empírica, o MPO \cite{vieira2022thesis} é apresentado predominantemente em linguagem natural, o que implica a ausência de uma representação formal que explicite, de maneira rigorosa, seus conceitos, relações e pressupostos semânticos.

A falta de representação formal torna o MPO fortemente dependente de interpretações humanas, abrindo espaço para ambiguidades semânticas, sobreposições conceituais e inconsistências na compreensão e aplicação de seus elementos. Como consequência, diferentes iniciativas que afirmam adotar o MPO podem fazê-lo de maneira semanticamente incompatível, ainda que mantenham aderência superficial à terminologia do modelo. Por exemplo, o termo \textit{indicador} pode ser interpretado por um observatório como uma métrica bruta extraída diretamente de uma fonte de dados, e por outro como o resultado de um processo de análise sobre múltiplas métricas --- ambos reportando-se ao mesmo termo do MPO, mas referindo-se a entidades conceitualmente distintas. Sem uma definição formal que explicite a natureza, os atributos e as relações de cada conceito, tais divergências permanecem implícitas e não verificáveis, comprometendo o uso do MPO como referência conceitual estável para análise, comparação e evolução de observatórios de projetos.

A OntoMPO endereça esse problema não por meio de validação empírica em observatórios reais --- o que está fora do escopo desta pesquisa ---, mas por fornecer uma
especificação de referência explícita e verificável: cada conceito do MPO passa a ter sua categoria ontológica, atributos e relações formalmente
definidos, de modo que iniciativas que pretendam alinhar-se ao modelo possam confrontar suas próprias interpretações com essa especificação. A contribuição da ontologia está, portanto, no nível da especificação
do modelo de referência, e sua validação com observatórios reais constitui
trabalho futuro.

\textbf{Assim, o problema de pesquisa desta dissertação consiste em como formalizar o Modelo para Observatórios de Projetos (MPO), de modo a explicitar seus conceitos, relações e restrições semânticas, reduzindo ambiguidades interpretativas e possibilitando sua representação formal e utilização em contextos computacionais e analíticos?}


\section{Proposta de solução}
Para enfrentar as limitações decorrentes da natureza predominantemente textual e informal do MPO \cite{vieira2022thesis}, esta pesquisa propõe a construção de uma ontologia de domínio que forneça uma formalização para o modelo. Essa ontologia tem como objetivo explicitar, de maneira sistemática, os conceitos, relações, atributos e pressupostos semânticos que compõem o MPO, tornando explícitos os compromissos ontológicos subjacentes ao modelo.

Outras abordagens formais, como DSLs ou perfis UML, formalizam restrições
estruturais, mas pressupõem conceitos já semanticamente definidos. O problema
do MPO é anterior a isso: a imprecisão semântica dos próprios conceitos.
A abordagem ontológica, ao impor compromissos categoriais explícitos, atua
diretamente nessa camada, além de favorecer a reutilização e a
interoperabilidade inerentes a um modelo de referência.

Por meio da formalização ontológica, busca-se reduzir ambiguidades conceituais, eliminar sobreposições semânticas e assegurar maior consistência interpretativa na compreensão e aplicação do MPO. A ontologia estabelece um vocabulário conceitual padronizado e semanticamente bem definido, capaz de servir como referência estável para a análise, comparação e evolução de iniciativas de observatórios de projetos fundamentadas no modelo.

Por fim, ressalta-se que a ontologia desenvolvida neste trabalho possui caráter essencialmente conceitual, tendo como objetivo principal a formalização do domínio do MPO. Dessa forma, não se pretende, neste momento, a sua aplicação direta em sistemas computacionais, mas sim o estabelecimento de uma base semântica rigorosa que possa apoiar futuras extensões e possíveis aplicações.

\section{Objetivos}
\subsection{Objetivo geral}
O objetivo geral desta pesquisa é desenvolver uma ontologia de referência para o MPO \cite{vieira2022thesis}, com o propósito de formalizar seus conceitos, relações, atributos e restrições semânticas, fornecendo uma fundamentação ontológica rigorosa que reduza ambiguidades conceituais, explicite os compromissos ontológicos do modelo e assegure maior consistência interpretativa em sua compreensão e aplicação.

Para alcançar esse objetivo, a ontologia é construída seguindo as etapas definidas pela metodologia \textit{Methontology} \cite{fernandez1997methontology}, que orienta de forma sistemática os processos de especificação, conceitualização, formalização, integração, implementação e avaliação de ontologias, assegurando rigor metodológico ao desenvolvimento da proposta.

\subsection{Objetivos específicos}

\begin{itemize}
    \item Realizar uma revisão exploratória da literatura sobre ontologias, fundamentação ontológica e observatórios de projetos, com o objetivo de identificar abordagens, conceitos e práticas relevantes ao domínio de estudo;
    \item Identificar, analisar e organizar os conceitos, relações e atributos presentes no MPO, considerando suas camadas, subcamadas e definições, a partir dos documentos oficiais do modelo e da literatura relacionada a observatórios de projetos;
    \item Estruturar o vocabulário conceitual do MPO por meio da elaboração dos artefatos previstos na metodologia \textit{Methontology}, incluindo glossário de termos, árvore de classificação de conceitos, dicionário de verbos e tabelas de regras e restrições, de modo a explicitar o significado e o papel de cada conceito no domínio;
    \item Formalizar os conceitos do MPO em um modelo ontológico, definindo classes, propriedades, axiomas, restrições e relações necessárias para representar o domínio de forma semanticamente consistente e formalmente precisa;
    \item Avaliar e validar a ontologia construída com base nas questões de competência, assegurando sua aderência ao escopo do MPO e seu potencial como base conceitual para futuras aplicações e investigações no domínio de observatórios de projetos.
\end{itemize}

\section{Delimitação da pesquisa}
Esta pesquisa está delimitada ao desenvolvimento de uma ontologia de referência para o MPO \cite{vieira2022thesis}, considerando exclusivamente os conceitos, relações, atributos e estruturas definidos na versão final do modelo proposta por \citeauthor{vieira2022thesis}. O estudo não tem como objetivo redefinir, estender ou modificar conceitualmente o MPO, mas sim formalizar ontologicamente seu conteúdo, tornando explícitos seus significados, relações e compromissos ontológicos conforme estabelecidos no modelo original.

No que se refere à sua classificação ontológica, a ontologia desenvolvida situa-se, quanto ao grau de formalismo, entre ontologias leves e ontologias pesadas, uma vez que apresenta uma estrutura conceitual rigorosa, com definições explícitas de classes, relações, propriedades e restrições, sem, contudo, empregar axiomatizações lógicas altamente expressivas ou mecanismos de raciocínio complexos. Em relação ao nível de generalidade, trata-se de uma ontologia intermediária, pois não se configura como uma ontologia de topo nem como uma ontologia de aplicação específica, atuando como elo conceitual entre uma ontologia fundamental e possíveis ontologias de aplicação no domínio de observatórios de projetos. Quanto à sua finalidade, a ontologia é classificada como uma ontologia de domínio, por formalizar conhecimentos específicos relacionados ao contexto de observatórios de projetos.

A construção da ontologia segue estritamente as etapas definidas pela metodologia \textit{Methontology} \cite{fernandez1997methontology}, restringindo-se aos artefatos, processos e níveis de formalidade previstos por essa metodologia. Não são adotadas outras metodologias ou \textit{frameworks} de desenvolvimento ontológico, tampouco realizadas comparações entre abordagens metodológicas distintas.

Por fim, esta pesquisa não abrange o desenvolvimento de sistemas computacionais. Sua implementação restringe-se à formalização da ontologia em uma linguagem adequada, com o apoio de ferramentas como o Protégé. O foco está na construção e na validação conceitual da ontologia proposta, não contemplando processos de validação empírica com usuários finais, organizações ou observatórios reais. As atividades de avaliação limitam-se à verificação da coerência interna do modelo e ao atendimento das questões de competência estabelecidas na etapa de especificação.

\section{Ambiente de pesquisa, publicações e outras atividades realizadas}
\subsection{Publicações}
\begin{itemize}
    \item \textbf{Observatório de Projetos de Extensão e Pesquisa: Uma perspectiva à luz de uma Universidade}. In: \textit{CISTI 2024} \cite{jose2024cisti}.
\end{itemize}

\section{Organização do documento}

Este documento está estruturado em cinco capítulos principais, além dos elementos pré-textuais e pós-textuais. A seguir, apresenta-se uma visão geral do conteúdo abordado em cada capítulo.

O \textbf{Capítulo 1} introduz o trabalho, apresentando a contextualização e a motivação da pesquisa, a descrição do problema, a proposta de solução, os objetivos e a delimitação do estudo, bem como o ambiente da pesquisa e as publicações e atividades associadas ao seu desenvolvimento.

O \textbf{Capítulo 2} corresponde à Fundamentação Teórica, na qual são discutidos os conceitos essenciais para a compreensão do trabalho. Esse capítulo aborda aspectos relacionados à observação e à transparência no contexto do gerenciamento de projetos, os fundamentos teóricos sobre observatórios, os conceitos do MPO \cite{vieira2022thesis}, bem como noções fundamentais sobre ontologias e metodologias para sua construção.

O \textbf{Capítulo 3} descreve a metodologia adotada para o desenvolvimento da ontologia proposta, com base na metodologia \textit{Methontology} \cite{fernandez1997methontology}.

O \textbf{Capítulo 4} apresenta os resultados obtidos e a respectiva discussão, contemplando os artefatos produzidos ao longo do processo de construção da ontologia. Nesse capítulo, discute-se como esses artefatos representam, organizam e formalizam o conhecimento do domínio do MPO.

Por fim, o \textbf{Capítulo 5} apresenta a conclusão do trabalho, destacando as principais contribuições da pesquisa, suas limitações e as oportunidades para trabalhos futuros relacionados ao desenvolvimento e à aplicação da ontologia proposta.
